%!TEX root = ../../main.tex
\section{한계}
\label{sec:disc-limitations}

\begin{enumerate}[leftmargin=2em]
  \item \textbf{탐색적 지위.} 유연한 기준선, CORP, 외부, 규모분리 패널은 시험 패치 15.16의 이전 노출 뒤에 생산되었다. 구간은 고정된 평가기와 선정 아래의 경기 부트스트랩 대비이지 사전 등록된 검정이나 모델 적합 불확실성이 아니다.
  \item \textbf{모델 정의 라벨.} 방향 라벨은 동결 평가기의 변화 부호이며 절대적 교전 가치는 관측되지 않는다. 모든 예측 가능성 주장은 그 평가기와 결과 시점 규칙에 조건부이다.
  \item \textbf{정의의 사람 검증 부재.} 독립적 경계 주석이 없으므로 실제 한타 탐지 정확도를 주장하지 않는다. 킬 없는 전투는 정의상 제외되고 이 텔레메트리로는 측정 불가능하다.
  \item \textbf{코호트 범위.} pick(한쪽 참여 1명 이하, 15.16 교전의 약 \sharePickTest\%)은 사전에 제외되며 이 평가기 아래의 결과가 없다. 두 코호트는 크기, 양성 비율, 균형 구간 비율이 달라 절대 점수를 비교하지 않는다.
  \item \textbf{보정기 변형.} 소규모 교전에서 두 단계 선택 규칙은 모든 모델에 시그모이드를 골랐을 것이다. 계약대로의 항등을 주로, 변형을 민감도로 보고하며 부호는 같다. 합동 대비의 결론은 이 선택에 의존한다.
  \item \textbf{외부 평가.} 점수 전용이고 구간이 없으며 어댑터가 없다. 외부 전 상태에서의 평가기 Brier는 외부 라벨을 스스로 검증하지 않는다.
  \item \textbf{정의 상수 민감도.} 보조 상수의 민감도는 선행 연구 시기나 파일럿 코퍼스에서 측정되었고 현재 코퍼스와 주 목표에서 재실행되지 않았다.
  \item \textbf{학습기와 정보 구성.} 하나의 로지스틱 사양을 사전에 고정했고, 트리 후보는 다른 입력을 썼으며, 정보군 모듈은 실행되지 않았다. 아키텍처 순위나 정보의 본질적 유무를 주장하지 않는다.
  \item \textbf{다중성.} 열네 개 안팎의 부트스트랩 구간에 대해 다중성 보정을 하지 않았다.
\end{enumerate}
